\documentclass[12pt]{article}

\usepackage[top=1in, bottom=1in, left=1in, right=1in]{geometry}
\usepackage{amsmath}
\usepackage{amssymb}
\usepackage{algorithm}
\usepackage{algpseudocode}
\usepackage{multirow}
\usepackage{mdframed}
\usepackage{booktabs}

% ── Bold algorithm keywords to match Dr. Alsaffar's style ───────────────
\algrenewcommand\algorithmicif{\textbf{if}}
\algrenewcommand\algorithmicelse{\textbf{else}}
\algrenewcommand\algorithmicend{\textbf{end}}
\algrenewcommand\algorithmicreturn{\textbf{return}}
\algrenewcommand\algorithmicwhile{\textbf{while}}
\algrenewcommand\algorithmicdo{\textbf{do}}
\algrenewcommand\algorithmicthen{\textbf{then}}
\algrenewcommand\algorithmicfor{\textbf{for}}
\algnewcommand{\algorithmicto}{\textbf{to}}

% ── Small-caps algorithm names ───────────────────────────────────────────
\newcommand{\algoname}[1]{\textsc{#1}}

% ── First line of algorithm = "Algorithm NAME(params)" like Dr. Alsaffar ─
\newcommand{\AlgorithmLine}[1]{\State \textbf{Algorithm} \algoname{#1}}

\begin{document}

% ── Header ───────────────────────────────────────────────────────────────
\begin{center}
    {\Large University of Bahrain}\\[4pt]
    College of Information Technology\\
    Department of Computer Science\\[6pt]
    ITCS347 Design and Analysis of Algorithms\\
    Second Semester 2025/2026\\[4pt]
    {\large Home Assignment \#2}
\end{center}

\bigskip

% ── Student information table ─────────────────────────────────────────────
\noindent
\begin{tabular}{|l|p{8cm}|p{1.8cm}|}
\hline
Student Name & Mohsen Abdulelah Ali / Abdulnaser Fawzi Hamadi & \multirow{3}{*}{\centering GRADE} \\
\cline{1-2}
Student ID   & 202305577 / 202008866 & \\
\cline{1-2}
Section      & 3 & \\
\hline
\end{tabular}

\bigskip\bigskip

% ── Problem statement ─────────────────────────────────────────────────────
\noindent\textbf{Civil Engineering -- Resultant Force of Components}\\[2pt]
When multiple forces act at different angles, the resultant force is:
\[
    F = \sqrt{\left(\sum F_x\right)^2 + \left(\sum F_y\right)^2}
\]
where $F_x$ and $F_y$ are horizontal and vertical components stored in two
arrays $Fx[0..n-1]$ and $Fy[0..n-1]$ of real numbers respectively.

\bigskip

% ═════════════════════════════════════════════════════════════════════════
% Question 1
% ═════════════════════════════════════════════════════════════════════════
\noindent\textbf{1.} Design a divide-and-conquer algorithm
\algoname{SumArray}$(A, left, right)$ that computes the sum of all elements
in a subarray $A[left..right]$.

\medskip

\begin{mdframed}
\noindent\textbf{Solution:}

\smallskip
\noindent
The algorithm divides $A[left..right]$ into two equal halves, recursively
computes the sum of each half, and combines the two partial sums by a single
addition (the combine step).

\medskip

\begin{algorithm}[H]
\begin{algorithmic}[1]
\AlgorithmLine{SumArray$(A,\ left,\ right)$}
\If{$left = right$}
    \State \textbf{return} $A[left]$
\Else
    \State $mid \leftarrow \lfloor(left + right) / 2\rfloor$
    \State $leftSum  \leftarrow \algoname{SumArray}(A,\ left,\ mid)$
    \State $rightSum \leftarrow \algoname{SumArray}(A,\ mid+1,\ right)$
    \State \textbf{return} $leftSum + rightSum$
\EndIf
\end{algorithmic}
\end{algorithm}
\end{mdframed}

\bigskip

% ═════════════════════════════════════════════════════════════════════════
% Question 2
% ═════════════════════════════════════════════════════════════════════════
\noindent\textbf{2.} Using \algoname{SumArray}, design an algorithm
\algoname{ComputeResult}$(Fx, Fy, n)$ that computes the resultant force $F$.

\medskip

\begin{mdframed}
\noindent\textbf{Solution:}

\smallskip
\noindent
\algoname{ComputeResult} calls \algoname{SumArray} on each component array to
obtain $\sum F_x$ and $\sum F_y$, then applies the formula directly.

\medskip

\begin{algorithm}[H]
\begin{algorithmic}[1]
\AlgorithmLine{ComputeResult$(Fx,\ Fy,\ n)$}
\State $sumX \leftarrow \algoname{SumArray}(Fx,\ 0,\ n-1)$
\State $sumY \leftarrow \algoname{SumArray}(Fy,\ 0,\ n-1)$
\State \textbf{return} $\sqrt{\,sumX^2 + sumY^2\,}$
\end{algorithmic}
\end{algorithm}
\end{mdframed}

\bigskip

% ═════════════════════════════════════════════════════════════════════════
% Question 3
% ═════════════════════════════════════════════════════════════════════════
\noindent\textbf{3.} Analyze the algorithm and find the time complexity
recurrence relation $T(n)$.

\medskip

\begin{mdframed}
\noindent\textbf{Solution:}

\smallskip
\begin{itemize}
    \item \textbf{Problem size:} $n = right - left + 1$
          (number of elements in the subarray).
    \item \textbf{Basic operation:} the addition in the combine step
          (\textbf{return} $leftSum + rightSum$).
\end{itemize}

\noindent\textbf{Analysis:}
\begin{itemize}
    \item \textbf{Base case} ($n = 1$, i.e., $left = right$):
          the algorithm returns $A[left]$ immediately; no addition is
          performed, so $T(1) = 0$.
    \item \textbf{Recursive case} ($n > 1$):
          two recursive calls are made, each on a subproblem of size
          $n/2$, and one addition is performed to combine the results.
\end{itemize}

\noindent The recurrence relation for \algoname{SumArray} is:
\[
    T(n) =
    \begin{cases}
        0              & \text{if } n = 1 \\[6pt]
        2\,T(n/2) + 1  & \text{if } n > 1
    \end{cases}
\]
\end{mdframed}

\bigskip

% ═════════════════════════════════════════════════════════════════════════
% Question 4
% ═════════════════════════════════════════════════════════════════════════
\noindent\textbf{4.} Prove that $T(n)$ is linear (for the algorithm
\algoname{SumArray} use the \textit{Master Theorem}).

\medskip

\begin{mdframed}
\noindent\textbf{Solution:}

\smallskip
\noindent
Applying the Master Theorem to $T(n) = a\,T(n/b) + \Theta(n^d)$, we
extract:
\[
    a = 2,\quad b = 2,\quad d = 0
    \qquad\bigl(\text{since } f(n) = 1 = n^0 \text{ is constant}\bigr)
\]

\noindent Compute $b^{\,d} = 2^0 = 1$.

\medskip
\noindent Since $a > b^{\,d}$ (i.e., $2 > 1$), we apply \textbf{Case 3}
of the Master Theorem:
\[
    T(n) \;\in\;
    \Theta\!\left(n^{\,\log_b a}\right)
    \;=\;
    \Theta\!\left(n^{\,\log_2 2}\right)
    \;=\;
    \Theta\!\left(n^{\,1}\right)
    \;=\;
    \Theta(n)
\]

\noindent Therefore $T(n) \in \Theta(n)$, which proves that
\algoname{SumArray} runs in \textbf{linear time}. \qquad $\blacksquare$
\end{mdframed}

\end{document}
